\section{Interval comparison and prospective trial use}

The efficiency bound translates into a comparison of asymptotic Wald interval lengths for a specified release and estimator. We use the known-assignment custom-Laplace construction of \citet{OhnishiAwan2025}, applied to the binary trial model in \cref{obj:def:binary-trial-model}. Its release \leanref{sym:\widetilde A_i}{\ensuremath{\widetilde A_i}} adds noise to an inverse-probability-weighted outcome contribution, and its estimator \leanref{sym:\widehat\tau_{\mathrm{OA}}}{\ensuremath{\widehat\tau_{\mathrm{OA}}}} averages these releases. The following definition fixes both components of the benchmark.

\begin{definitionv}{def:oa-release}[Custom-Laplace release $\widetilde A_i$]
The custom-Laplace release and its sample-mean estimator are
\[
\widetilde A_i
=
\frac{W_iY_i}{p}
-
\frac{(1-W_i)Y_i}{q}
+
L_i,
\qquad
\widehat\tau_{\mathrm{OA}}
=
\frac{1}{n}\sum_{i=1}^{n}\widetilde A_i.
\]
The noise variables \(L_i\) are independent and identically distributed, independent of all subject records, with
\[
L_i\sim\operatorname{Lap}(0,\Delta_A/\varepsilon),
\qquad
\Delta_A=p^{-1}+q^{-1}=\frac{1}{pq}.
\]
\end{definitionv}

The weighted outcome contribution has expectation equal to the treatment effect under randomized assignment, following the inverse-probability-weighting principle of \citet{Horvitz1952}. The independent, centered noise \leanref{sym:L_i}{\ensuremath{L_i}} preserves this expectation. Its scale is determined by the sensitivity constant \leanref{sym:\Delta_A}{\ensuremath{\Delta_A}}, the distance between the largest and smallest weighted contributions on the four-record alphabet. Thus the benchmark combines variation from assignment and outcomes with the variance of the added Laplace noise.

Write \leanref{sym:V_{\mathrm{OA}}(\theta,p,\varepsilon)}{\ensuremath{V_{\mathrm{OA}}(\theta,p,\varepsilon)}} for the asymptotic variance constant of this sample mean. Comparing it with \(V^*(\theta,p,\varepsilon)\), defined in \cref{obj:def:staircase-saddle}, gives an interval-length comparison at a common sample size and coverage level. The next result supplies the benchmark variance, its consistent estimator, and the corresponding comparison with the selected private-pilot procedure.

\begin{theoremv}{thm:laplace-comparison}[Laplace variance and interval comparison]
Let \(\theta=(\mu_0,\mu_1)^\top\) be the pair of arm means, \(p\) the treatment-assignment probability, \(\varepsilon\) the privacy parameter, and \(\gamma\) the Wald error probability. Suppose the following conditions hold, under the usual measurability conditions.
\begin{itemize}
\item \textbf{(Binary randomized trial.)} The subjects \(\{(Y_i(0),Y_i(1),W_i):i\ge1\}\) are independent and identically distributed. The potential outcomes are binary, treatment assignment is independent of the potential outcomes, and
\[
\Pr(W_i=1)=p,\qquad
Y_i=(1-W_i)Y_i(0)+W_iY_i(1),\qquad
\mu_0=\mathbb E[Y_i(0)],\qquad
\mu_1=\mathbb E[Y_i(1)],
\]
with \(0<p<1\) and \(0<\mu_0,\mu_1<1\). Set \(q=1-p\) and \(\tau=\mu_1-\mu_0\).
\item \textbf{(Public experiment and representation.)} The specified public experiment is \(P_\theta^{\otimes n}\) for every positive sample size \(n\), and the private records \(X_i=(W_i,Y_i)\), ordered as \((0,0),(0,1),(1,0),(1,1)\), have this joint law. Their one-subject probabilities are
\[
\bigl(q(1-\mu_0),q\mu_0,p(1-\mu_1),p\mu_1\bigr).
\]
\item \textbf{(Comparator noise.)} The pairs \((X_i,L_i)\) are jointly independent across subjects, each \(L_i\) is independent of \(X_i\), and
\[
L_i\sim\operatorname{Lap}(0,\Delta_A/\varepsilon),
\qquad
\Delta_A=\frac1p+\frac1q.
\]
\item \textbf{(Privacy and critical value.)} We have \(\varepsilon>0\), \(0<\gamma<1\), and
\[
z_{1-\gamma/2}
=\inf\{z\in\mathbb R:\Phi(z)\ge1-\gamma/2\}>0,
\]
where \(\Phi\) is the standard-normal distribution function.
\item \textbf{(Strong-saddle selection.)} Fix any measurable strong-saddle selector
\(\vartheta\mapsto(\alpha_\vartheta,t_\vartheta,J_\vartheta)\)
as in \cref{obj:def:pilot-estimator}. In particular, at every interior \(\vartheta\), its selected weights maximize the profiled objective at its selected minimizing coordinate:
\[
F_\vartheta(\alpha,t_\vartheta)
\le F_\vartheta(\alpha_\vartheta,t_\vartheta)
\qquad\text{for every }\alpha\in\mathcal A_\varepsilon.
\]
\end{itemize}

Define the comparator releases, their sample mean, and their empirical variance by
\[
\widetilde A_i
=\frac{W_iY_i}{p}-\frac{(1-W_i)Y_i}{q}+L_i,
\qquad
\widehat\tau_{\mathrm{OA}}
=\frac1n\sum_{i=1}^n\widetilde A_i,
\qquad
\widehat V_{\mathrm{OA}}
=\frac1n\sum_{i=1}^n\widetilde A_i^2
-\widehat\tau_{\mathrm{OA}}^2,
\]
and set
\[
V_{\mathrm{OA}}(\theta,p,\varepsilon)
=\frac{\mu_1}{p}+\frac{\mu_0}{q}-\tau^2
+\frac{2\Delta_A^2}{\varepsilon^2}.
\]
The specified public experiment and its private-record representation hold jointly with the following conclusions:
\[
\sqrt n\bigl(\widehat\tau_{\mathrm{OA}}-\tau\bigr)
\rightsquigarrow
N\!\left(0,V_{\mathrm{OA}}(\theta,p,\varepsilon)\right),
\]
\[
\frac{z_{1-\gamma/2}\sqrt{V_{\mathrm{OA}}(\theta,p,\varepsilon)/n}}
     {z_{1-\gamma/2}\sqrt{V^*(\theta,p,\varepsilon)/n}}
\longrightarrow
\sqrt{\frac{V_{\mathrm{OA}}(\theta,p,\varepsilon)}
           {V^*(\theta,p,\varepsilon)}}.
\]
Moreover, for every \(\delta>0\),
\[
\Pr\!\left(
\left|\widehat V_{\mathrm{OA}}
      -V_{\mathrm{OA}}(\theta,p,\varepsilon)\right|>\delta
\right)\longrightarrow0.
\]

There exists a pilot schedule satisfying
\cref{obj:ass:pilot-diverges,obj:ass:pilot-sublinear}, witnessed by
\[
m_n=\lfloor\sqrt n\rfloor.
\]
For this schedule, use the selected pilot procedure in
\cref{obj:def:pilot-estimator} and its variance statistic \(\widehat V^*\).
For every sequence of couplings whose first marginal is the trial probability law and whose second marginal is the transcript law of that pilot procedure at \((\theta,p)\), and for every \(\delta>0\), the coupled probabilities satisfy
\[
\Pr\!\left(
\left|
\sqrt{\frac{\widehat V_{\mathrm{OA}}}{\widehat V^*}}
-\sqrt{\frac{V_{\mathrm{OA}}(\theta,p,\varepsilon)}
             {V^*(\theta,p,\varepsilon)}}
\right|>\delta
\right)\longrightarrow0.
\]

At \((p,\mu_0,\mu_1,\varepsilon)=(1/2,1/2,1/2,1)\),
\[
V_{\mathrm{OA}}=34,
\qquad
V^*=\left(\frac{\cosh(1/2)}{\sinh(1/2)}\right)^2
< V_{\mathrm{OA}}.
\]
\end{theoremv}

The variance decomposition explains the comparison. The first three terms of \(V_{\mathrm{OA}}\) describe the variability of the weighted binary outcome; the final term is the independent noise variance. At the stated balanced point, these contributions are \(2\) and \(32\), respectively. The optimized variance instead reflects the release chosen for efficient estimation of the treatment-effect contrast. A common normal critical value \leanref{sym:z_{1-\gamma/2}}{\ensuremath{z_{1-\gamma/2}}} and sample size cancel from the half-length ratio, leaving the square root of the variance ratio. The strict inequality at the displayed point therefore establishes shorter asymptotic intervals for the selected efficient procedure relative to the custom-Laplace sample mean.

The statistical comparison is specific to the release and sample-mean estimator in \cref{obj:def:oa-release}. Consistency of \(\widehat V_{\mathrm{OA}}\), together with the selected pilot variance statistic, makes the population ratio operational through estimated Wald half-lengths. The coupling conclusion permits any dependence between the two procedures consistent with their stated marginal laws. Its pilot schedule supplies both an increasing number of pilot observations and a vanishing pilot fraction, as required by \cref{obj:ass:pilot-diverges,obj:ass:pilot-sublinear}.

\subsection{Prospective attendance-threshold design}

The Colombian cash-transfer program studied by \citet{BarreraOsorioEtAl2011} motivates a prospective binary endpoint based on meeting its \(80\%\) attendance threshold. For a prospective randomized evaluation, let \(D_i(a)\) denote subject \(i\)'s fraction of scheduled school days attended under assignment \(a\), measured over a prespecified observation window. Define the proposed potential outcome by
\[
Y_i(a)=\mathbf 1\{D_i(a)\ge 0.80\},
\qquad a\in\{0,1\}.
\]
The observation window, denominator of scheduled days, and rules for recording attendance would be fixed before assignment. Under consistency, the observed indicator is \(Y_i=Y_i(W_i)\), and the treatment effect is the difference between the two probabilities of meeting the threshold. This construction uses the program's threshold to define an endpoint for a prospective trial.

The design inputs are the known assignment probability \(p\), the fixed privacy budget \(\varepsilon\), the planned sample size \(n\), the Wald error probability \(\gamma\), and interior planning values for the threshold probabilities \(\mu_0\) and \(\mu_1\). For independently sampled subjects with randomized assignment and the record representation in \cref{obj:def:binary-trial-model}, these inputs determine the two variance constants in \cref{obj:thm:laplace-comparison}. Implementation of the efficient procedure additionally uses the measurable strong-saddle selector in \cref{obj:def:pilot-estimator} and the stated pilot schedule. The planning arm means specify scenarios for evaluating expected precision; the private pilot supplies the parameter estimates used to select the release during the trial.

For illustration, take the hypothetical planning values
\((p,\mu_0,\mu_1,\varepsilon)=(1/2,1/2,1/2,1)\).
The exact witness in \cref{obj:thm:laplace-comparison} then gives the following prospective design calculations:
\[
\begin{aligned}
\text{Asymptotic half-length ratio at equal sample size}
&=\sqrt{34}\,\tanh(1/2)
\approx 2.69,\\
\text{Asymptotic sample-size ratio at equal half-length}
&=34\,\tanh^2(1/2)
\approx 7.26.
\end{aligned}
\]
Both ratios place the custom-Laplace sample mean in the numerator and the selected efficient procedure in the denominator. The second calculation follows by equating their asymptotic Wald half-lengths at a common coverage level. These are prospective calculations for the specified binary endpoint and planning values, translating the variance comparison into precision and recruitment quantities for trial design.
